% Generated draft: requires mathematical and source review.
\documentclass[11pt]{article}
\usepackage[margin=25mm]{geometry}
\usepackage{fontspec}
\setmainfont{Latin Modern Roman}
\usepackage{amsmath,amssymb,mathtools,unicode-math}
\setmathfont{Latin Modern Math}
\usepackage{enumitem,needspace,xurl,hyperref}
\setlength{\parindent}{0pt}
\setlength{\parskip}{5pt}
\setlength{\emergencystretch}{3em}
\newcommand{\N}{\mathbb{N}}
\newcommand{\Z}{\mathbb{Z}}
\newcommand{\Q}{\mathbb{Q}}
\newcommand{\R}{\mathbb{R}}
\newcommand{\C}{\mathbb{C}}
\newcommand{\F}{\mathbb{F}}
\newcommand{\A}{\mathbb{A}}
\newcommand{\PP}{\mathbb{P}}
\newcommand{\E}{\mathbb{E}}
\newcommand{\eps}{\varepsilon}
\newcommand{\dd}{\,\mathrm{d}}
\begin{document}
\section*{UnsolvedMath ID: 30005507}
\textbf{Source problem code:} catalogue entry
\subsection{Definitions and mathematical statement}
Let \(G\) be a finite group. Write \(\operatorname{Irr}(G)\) for its irreducible complex characters, and write \(\operatorname{Irr}(B)\) for the ordinary irreducible characters belonging to a \(2\)-block \(B\) of \(G\). A block is real if \(\operatorname{Irr}(B)\) is closed under complex conjugation, and nonprincipal if it does not contain the trivial character. The standard nilpotence condition is that \(N_G(Q,b_Q)/C_G(Q)\) is a \(2\)-group for every \(B\)-Brauer pair \((Q,b_Q)\), where \(b_Q\) is a block of the centralizer \(C_G(Q)\) and \(N_G(Q,b_Q)\) is its stabilizer in \(N_G(Q)\).

Let \((D,E)\) be a defect pair for \(B\), in the block-theoretic sense intended by the cited source, with \(D\) a defect group of \(B\). This is a substantive source-dependent hypothesis: the supplied record does not define the extended defect group \(E\) or the construction of a defect pair. In particular, the displayed character sum requires \(D\leq E\) and \(e^2\in D\) for every \(e\in E\setminus D\). These domain conditions alone are not a definition of a defect pair and must not replace that hypothesis.

For a positive integer \(n\), let \(v_2(n)\) denote its \(2\)-adic valuation. Put \(a=v_2(|G|)\) and \(d=v_2(|D|)\). The height of \(\chi\in\operatorname{Irr}(B)\) is
\[
 h_B(\chi)=v_2(\chi(1))-a+d.
\]
For \(\lambda\in\operatorname{Irr}(D)\), its height in the unique \(2\)-block of the \(2\)-group \(D\) is \(h_D(\lambda)=v_2(\lambda(1))\). The Frobenius--Schur indicator of \(\chi\in\operatorname{Irr}(G)\) is
\[
 \epsilon(\chi)=\frac{1}{|G|}\sum_{g\in G}\chi(g^2).
\]

\textbf{Conjecture 2, as supplied.} For every finite group \(G\), every real, nilpotent, nonprincipal \(2\)-block \(B\) of \(G\), and every defect pair \((D,E)\) covered by the source's convention, there exists a bijection
\[
 \Gamma:\operatorname{Irr}(D)\longrightarrow\operatorname{Irr}(B)
\]
such that, for every \(\lambda\in\operatorname{Irr}(D)\),
\[
 h_B(\Gamma(\lambda))=h_D(\lambda),
 \qquad
 \epsilon(\Gamma(\lambda))
 =\frac{1}{|D|}\sum_{e\in E\setminus D}\lambda(e^2).
\]
The precise defect-pair convention remains to be confirmed from the cited source before this draft can be regarded as a fully specified formulation.
\subsection{Short English statement}
The conjecture asks for a matching between the irreducible characters of a defect group and those in a real nilpotent nonprincipal \(2\)-block. The matching must preserve character heights and recover every matched character's Frobenius--Schur indicator from the stated sum over the extended defect group. The supplied record leaves the precise definition of a defect pair implicit.
\subsection{Sources}
\begin{itemize}
\item Supplied local UnsolvedMath registry record, permanent ID \(30005507\), titled ``Frobenius--Schur Indicators in Real Nilpotent Blocks,'' exact source problem code \texttt{catalogue entry}; Oberwolfach Reports Open Problems collection, proposed year \(2023\). No individual proposer is specified.
\item \emph{Representations of Finite Groups} (2023), Oberwolfach Reports, report 19/2023, pp.\ 1031--1084, DOI \texttt{10.4171/owr/2023/19}: \url{https://doi.org/10.4171/owr/2023/19}. The supplied original excerpt labels the statement ``Conjecture 2.'' The registry says its reconstructed statement was checked against the official report and labels it \texttt{corrected\_verified}; that is supplied source-verification metadata, not an independent verification performed for this draft.
\item Sambale, \emph{Real characters in nilpotent blocks}: \url{https://doi.org/10.1007/s10013-023-00623-5}.
\item Sambale, primary preprint: \url{https://arxiv.org/abs/2301.13440}.
\item Supplied literature assessment, checked \texttt{2026-08-22}: the record labels the problem partially solved and reports that Sambale's bijection and indicator formula hold when the defect is abelian or dihedral and for solvable or quasisimple groups, while the general conjecture remains unproved. This is supplied dated curation derived from the Oberwolfach curation, not an independently verified literature assessment. The record explicitly requests independent verification of time-sensitive status and cited claims.
\end{itemize}
\end{document}
